\section{Results}
\label{sec:results}

Results test whether partner-local precision tracks predictability over payoff, acts through the \(\beta_k \rightarrow \gamma_k\) deployment pathway, and remains behaviorally consequential under partner choice, betrayal, and profile variation. Unless otherwise noted, uncertainty is reported as a two-sided 95\% percentile bootstrap confidence interval over simulation seeds; matched conditions are resampled as paired seed differences, and correlations use a seed-cluster bootstrap.

\subsection{Partner-local precision tracks predictability rather than realized payoff}
\label{sec:predictability_value}

Partner-local precision was more strongly associated with partner-response likelihood surprisal than with realized payoff. Because the affective-precision update is itself driven by partner-response surprisal rather than payoff, we expected this relationship and use the full multi-partner simulation here to verify that the intended selectivity is preserved once confidence updating, partner choice, and policy selection interact. Because partner-choice episodes couple payoff, surprisal, and exposure, we computed active-encounter partial correlations controlling for encounter count and the alternative signal. The precision--surprisal association remained substantially stronger ($r=-0.660$, 95\% CI [$-0.766$, $-0.492$]) than the precision--payoff association ($r=0.094$, 95\% CI [$-0.112$, $0.291$]). Negative precision--surprisal association indicates that less surprising partner responses are associated with greater $\beta_k$-derived policy precision; Figure~\ref{fig:model_fitness_contrast} plots absolute magnitudes for readability.

\begin{figure}[H]
\centering
\includegraphics[width=\linewidth]{fig_model_fitness_beta_reward_divergence.pdf}
\caption{Partner-local precision is more strongly associated with partner-response likelihood surprisal than with realized payoff. Correlation panels plot absolute active-encounter partial correlations with 95\% seed-cluster bootstrap confidence intervals; the payoff panel shows mean payoff with 95\% bootstrap confidence intervals in the same predictability--payoff analysis run. Because partner-response surprisal is the designated evidence source for the precision update, this comparison verifies that the intended selectivity remains present in the full multi-partner simulation. Pooling precision evidence across partners attenuates the relationship-local predictability signature.}
\label{fig:model_fitness_contrast}
\end{figure}

Pooling precision evidence across partners weakened this relationship-specific signal. In the shared-$\beta$ ablation, per-partner beliefs remained separate but the precision tracker was updated from pooled partner evidence. The precision--surprisal association was weaker ($r=-0.454$, 95\% CI [$-0.619$, $-0.275$]), while the precision--payoff association remained small ($r=0.072$, 95\% CI [$-0.163$, $0.390$]). Thus, pooling preserves separate partner beliefs but weakens the alignment between confidence and the relationship-specific predictive history from which that confidence is derived.

Payoff in the same predictability--payoff analysis run provides behavioral context rather than an independent definition of the mechanism. Although payoff varies across runtime variants, the partner-local precision signal remains much more strongly associated with partner-response predictability than with realized payoff. The payoff panel therefore complements the correlation panels without changing the central interpretation: partner-local affective precision tracks confidence in the current partner model rather than directly tracking reward.

\subsection{Tracked-only ablation isolates policy commitment}
\label{sec:deployment}

To test whether partner-local affective precision changes behavior through policy deployment rather than partner-state inference, we used a tracked-only ablation. This variant updates $q(\beta_k)$ normally while holding $\gamma_k=\gamma_{\text{base}}$, preserving the precision estimate but preventing it from modulating policy selection.

The tracked-only ablation confirmed that the tracker can move without producing the same deployment effect. In the open graded deployment run, partner-local and tracked-only conditions showed comparable within-episode variation in posterior mean $\bar{\beta}_k$, differing by only $0.043$ in mean range. The key difference was deployment: full precision modulation lowered policy entropy by $0.770$ nats relative to tracked-only (95\% CI [$-0.966$, $-0.577$]). Cumulative payoff was also somewhat higher under partner-local precision ($2003.3$ versus $1966.8$; difference $=36.5$, 95\% CI [$-8.6$, $84.1$]). Thus, $\beta_k$ matters behaviorally when it is allowed to modulate $\gamma_k$, producing sharper policy commitment. When that pathway is cut, the tracker remains active but the same entropy reduction is lost.

\begin{figure}[H]
\centering
\includegraphics[width=\linewidth]{fig_deployment_social_summary.pdf}
\caption{Open graded regime: tracked-only and partner-local precision show comparable movement in posterior mean $\bar{\beta}_k$, but only partner-local precision deploys that tracker through $\gamma_k$. Relative to tracked-only, partner-local precision lowers policy entropy, isolating the sharpening effect to the $\beta_k \rightarrow \gamma_k$ policy-precision pathway. Bars and whiskers show seed means or matched-seed differences with 95\% bootstrap confidence intervals; payoff is included as a secondary behavioral outcome.}
\label{fig:deployment_contrast}
\end{figure}

This supports the interpretation of affective precision as a calibration layer rather than an inference-improvement mechanism. The tracker does not directly change the partner-state update; it regulates how strongly the agent commits to policies based on its current partner model. Whether sharper commitment improves payoff depends on the social context, not on precision alone.

\subsection{Partner choice extends policy commitment across social options}
\label{sec:partner_choice}

Using the same underlying partner-choice trajectories analyzed in Section~\ref{sec:deployment}, we next examine a distinct readout: how precision deployment shapes commitment across partner--investment policies and how selected interactions are distributed across partner types. In this regime, each candidate pairs a partner with a multi-step investment plan, of which only the first step is executed each round. Precision modulation lowered policy entropy relative to the no-affect control ($6.887$ versus $7.657$; difference $=-0.770$, 95\% CI [$-0.966$, $-0.577$]), indicating that the $\beta_k \rightarrow \gamma_k$ pathway shapes commitment over partner--investment policies, not only investment levels within a fixed interaction.

Partner selection itself remained broadly distributed across the four partner types (cooperator $29.6\%$ versus $30.7\%$; exploiter $24.4\%$ versus $22.6\%$; reciprocator $22.0\%$ versus $24.0\%$; random $24.0\%$ versus $22.7\%$). The type-specific differences were small, and the paired 95\% bootstrap interval for each selected-type difference included zero. These allocation percentages therefore should not be read as evidence of a general preference for any one partner type. The main effect remains sharper commitment over partner--investment policies rather than systematic selection of a particular social type.

Together, these results show that partner-local precision can alter policy commitment without imposing a fixed social preference. Whether sharper commitment helps or hurts depends on whether the social environment remains stable. The abrupt-betrayal condition tests this boundary directly.

\subsection{Abrupt betrayal tests confidence under change}
\label{sec:betrayal}

Accumulated confidence helps when it remains attached to a valid partner model, but can lag behind abrupt social change. The abrupt-betrayal condition tests this boundary: partner~0 begins as an exploiter in a trusting stance and switches to a hostile stance at round~31, and the agent must adapt while carrying confidence built from the pre-switch relationship.

Precision modulation continued to affect deployment after the switch. Mean policy entropy was lower under partner-local precision than under the no-affect control, with a mean difference of $-2.149$ nats (95\% CI [$-2.379$, $-1.895$]). Joint type--stance accuracy was also higher by $0.158$ (95\% CI [$0.084$, $0.226$]), consistent with altered engagement changing the evidence available to the agent. Because the tracked-only ablation in Section~\ref{sec:deployment} isolated the $\beta_k \rightarrow \gamma_k$ pathway, this accuracy difference is best interpreted as a downstream effect of changed engagement and sampling, not as a direct improvement to the partner-state update.

\begin{figure}[H]
\centering
\includegraphics[width=\linewidth]{fig_betrayal_boundary_summary.pdf}
\caption{Abrupt betrayal in the partner-choice regime. The dashed line marks the round-31 stance switch for partner~0. Shaded bands show 95\% seed-bootstrap confidence intervals. Partner-local precision maintains lower policy entropy after the switch, while the tracked-only condition shows that $\beta_k$ dynamics can persist without the same deployment effect. Betrayed-partner selection illustrates the reallocation problem: confidence remains behaviorally relevant after the relationship changes, but downstream outcomes depend on how engagement is redistributed.}
\label{fig:betrayal_boundary}
\end{figure}

Payoff was also higher in this regime, with a cumulative difference of $58.2$ (95\% CI [$37.5$, $82.3$]). This is consistent with a calibration mechanism: confidence can keep policies decisive after a relationship changes, but whether that commitment improves payoff depends on how subsequent engagement is redistributed. The payoff result is therefore specific to the tested social environment rather than evidence that greater precision generally maximizes reward.

The betrayal result exposes a timing problem in social confidence: confidence built from reliable past evidence can remain behaviorally active after the relationship changes. We next vary precision gain and prior model fitness to examine how confidence-revision dynamics differ across agents.

\FloatBarrier
\subsection{Precision dynamics as computational profiles of social trust calibration}
\label{sec:phenotypes}

The betrayal result shows that confidence revision has temporal structure: confidence can be too weak, too slow, or too reactive. Varying precision gain $\alpha$ confirmed that gain controls the amplitude of $\beta_k$ dynamics: in the betrayal regime, the average within-episode max--min range of posterior mean $\bar{\beta}_k$ increased monotonically from $0.097$ at $\alpha=0.05$ to $0.675$ at $\alpha=8.0$ (Spearman $\rho=1.0$). Payoff did not follow this monotonic pattern, indicating that stronger confidence responsiveness is not automatically better calibration.

Crossing prior model fitness with gain produced distinct trust-calibration profiles. Low gain muted confidence dynamics whereas high gain amplified revision, while naive and cautious priors placed the agent at different initial levels of confidence in its partner models. These differences shaped how confidence accumulated and responded to subsequent social evidence. In the tested betrayal profiles, the anxious-reactive configuration (naive prior, high gain) achieved the highest mean payoff ($2285.8$), replacing the default configuration as the highest-payoff profile in the corrected model. The detailed payoff ordering therefore changed, but the profiles continued to express distinct confidence-revision dynamics. Exact priors, gain values, and updated profile metrics are reported in Appendix~\ref{app:extended_results}.

Trust-repair diagnostics separately examined reengagement, restored confidence, and payoff recovery after a partner returned to cooperative behavior. These readouts remained dissociable across parameter settings, indicating that renewed interaction need not coincide with identical confidence dynamics. Extended profile and repair metrics are reported in Appendix~\ref{app:extended_results}.
